% =====================================================================
% supplemental.tex
% Supplemental Material for the IEEE TIFS submission
% "Decision-Grade Masking: A Sound-by-Construction Static Lint
%  and Equipment-Free Verification Flow for Robust d-Probing
%  Security of a First-Order AES S-box".
%
% This file is the supplemental appendix material.  It contains
% proofs, witnesses, and experimental details that exceed the
% main paper's 12-page limit and that the reviewer can read in
% parallel with the main manuscript.pdf.  All claims are
% numbered to match the main paper.
%
% Build:  pdflatex supplemental.tex ; bibtex supplemental ;
%         pdflatex supplemental.tex ; pdflatex supplemental.tex
% (MiKTeX 25.12 + sttools package; 4-pass).
% =====================================================================

\documentclass[journal,onecolumn]{IEEEtran}

\usepackage{amsmath,amssymb,amsfonts}
\usepackage{amsthm}
\usepackage{graphicx}
\usepackage{booktabs}
\usepackage{multirow}
\usepackage{url}
\usepackage[hidelinks]{hyperref}
\usepackage{algorithm}
% Use the legacy 'algorithmic' package (same as main paper) so that
% \\REQUIRE / \\ENSURE / \\STATE / \\RETURN are consistent across both
% documents.  algpseudocode is *not* loaded here.
\usepackage{algorithmic}
\usepackage{listings}
\usepackage{xcolor}
\usepackage{tikz}
\usetikzlibrary{positioning,calc,shapes,arrows.meta}

% Allow long identifiers to break more permissively.
\sloppy
\tolerance=2000
\emergencystretch=15pt

\begin{document}
% Custom macros copied from the main paper preamble
% (kept identical so both documents share theorem environments and
% the GF/share/Adv/prob macros verbatim).
\newcommand{\GF}{\mathrm{GF}}
\newcommand{\share}[2]{#1^{(#2)}}
\newcommand{\Adv}{\mathcal{A}}
\newcommand{\prob}[1]{\Pr\!\left[#1\right]}

% Theorem-like environments (must match main paper for consistent
% numbering / cross-references)
\newtheorem{theorem}{Theorem}
\newtheorem{lemma}{Lemma}
\newtheorem{definition}{Definition}
\newtheorem{corollary}{Corollary}
\newtheorem{remark}{Remark}


\title{Supplemental Material for ``Probing-Resistant Masked Hardware: Formal Security Proofs and Simulation-Based Side-Channel Verification''\\
\large Proofs, Witnesses, and Experimental Details}

\author{Cheonshik Kim}

\markboth{Supplemental Material}{IEEE TIFS submission, 2026}

\maketitle

\begin{abstract}
This supplemental material contains (A) the corrected PROLEAD
worst-case $-\log_{10}(p)$ and the discrepancy discussion,
(B) the explicit witness for the gap
(Theorem~4 (main paper)) of the main paper, and (C) the pre-barrier PROLEAD history and the scope of the second-order instance.  All statements are numbered to match
the main paper; references are co-located in this archive's
\texttt{manuscript.bib}.
\end{abstract}

\renewcommand{\thesection}{S\arabic{section}}
\renewcommand{\thetheorem}{S\arabic{theorem}}
\renewcommand{\thelemma}{S\arabic{lemma}}
\renewcommand{\thecorollary}{S\arabic{corollary}}
\renewcommand{\theremark}{S\arabic{remark}}
\renewcommand{\thedefinition}{S\arabic{definition}}

\section{PROLEAD worst-case $-\log_{10}(p)$ correction}
\label{app:prolead-corr}

We record here the corrected PROLEAD worst-case
$-\log_{10}(p)$ for the first-order baseline, and explain the
discrepancy with an earlier artefact in the repository
(\texttt{syn/\allowbreak{}prolead/\allowbreak{}run\_\allowbreak{}full\_\allowbreak{}result.json}) that a reviewer
might otherwise re-discover.

The PROLEAD run reported in
the lint cross-check table (main paper, first row) was
re-issued against
\texttt{syn/\allowbreak{}prolead/\allowbreak{}circuit\_\allowbreak{}struct.v} with
\texttt{syn/\allowbreak{}prolead/\allowbreak{}config.json} ($512{,}000$ simulations,
$d = 1$).  The run is verdicted OKAY with $0 / 14{,}400$ probing
sets flagged, and the worst-case
$-\log_{10}(p)$ reported by PROLEAD is

\[
-\log_{10}(p_{\max}) \;=\; 3.386122.
\]

This value is the one quoted in the body of
Section~5 (main paper) and in the manifest
(``$3.386122$'').  The earlier artefact
\texttt{syn/\allowbreak{}prolead/\allowbreak{}run\_\allowbreak{}full\_\allowbreak{}result.json} reports two
inconsistent numbers, \texttt{worst\_\allowbreak{}log10p\_\allowbreak{}stdout:\allowbreak 4.205497}
and \texttt{max\_\allowbreak{}log10p:\allowbreak 3.3861216}, and its
\texttt{stdout\_\allowbreak{}log} field points at a different project
directory.  The first of these ($4.205497$) is the
$-\log_{10}(p)$ of a single intermediate
permutation-test iteration that was captured by accident; the
second ($3.3861216$) is the genuine worst case for the run and
agrees with the corrected value $3.386122$ to the precision
carried by PROLEAD.  We therefore use $3.386122$ throughout,
and we have retained the inconsistent artefact under
\texttt{legacy/\allowbreak} in the artifact repository only as a historical
record; it is excluded from the artifact manifest.

At $d = 1$ the Bonferroni-corrected threshold
$-\log_{10}(0.05 / 14{,}400) = 5.4594$ is essentially equal to
PROLEAD's fixed threshold of $5$, which is why the fixed
threshold is sound at first order.  The empirical evaluation in this paper is first-order only.

\section{Explicit Witness for Theorem~4 (main paper)}
\label{app:witness}

We give a fully worked algebraic construction demonstrating the gap
between value-based verification (which inspects only
register-output distributions) and cone-based lint
(Corollary~3 (main paper)). We first show that a glitch-extended probe
recovers the secret from the refresh-free first multiplication
$x^{2}\!\cdot\!x^{4}$ of the Itoh--Tsujii chain (Section~4 (main
paper)), and then give a minimal netlist whose registered value is
independent of the secret while the same recovery holds.

\subsection{Netlist under consideration}

Let $x\in\GF(2^{8})$ be a first-order additively shared secret,
$x = x^{(0)}\oplus x^{(1)}$. Because squaring is $\mathbb{F}_{2}$-linear
in characteristic $2$, the intermediate powers inherit a
canonical linear encoding:
\begin{equation}
(x^{2})^{(a)} = \bigl(x^{(a)}\bigr)^{2},
\qquad
(x^{4})^{(a)} = \bigl(x^{(a)}\bigr)^{4},
\qquad a\in\{0,1\}.
\label{eq:witness-linear-shares}
\end{equation}
A refresh-free DOM-style first cross-product of
$x^{2}\!\cdot\!x^{4}$ writes
\begin{equation}
z \;=\; (x^{2})^{(0)}\!\cdot\!(x^{4})^{(0)} \,\oplus\, c_{01} \,\oplus\, c_{10} \,\oplus\, (x^{2})^{(1)}\!\cdot\!(x^{4})^{(1)},
\label{eq:witness-z}
\end{equation}
with cross-domain terms
$c_{01} = (x^{2})^{(0)}\!\cdot\!(x^{4})^{(1)}$ and
$c_{10} = (x^{2})^{(1)}\!\cdot\!(x^{4})^{(0)}$. Let
$\mathrm{Cone}(c_{01})$ denote the glitch-extended combinational
fan-in feeding the register that latches $c_{01}$ at the cycle
boundary.

\subsection{Glitch-extended probe deterministically recovers $x$}

\begin{lemma}
\label{lem:witness-recover}
A single glitch-extended probe on $\mathrm{Cone}(c_{01})$ observes
$(x^{2})^{(0)}$ and $(x^{4})^{(1)}$; from these values, an offline
adversary computes $x$ deterministically. Consequently the pair
$\bigl((x^{2})^{(0)},(x^{4})^{(1)}\bigr)$ carries the full secret.
\end{lemma}

\begin{proof}
By construction of
$\mathrm{Cone}(c_{01})$ the two operands $(x^{2})^{(0)}$ and
$(x^{4})^{(1)}$ are combinationally driven into the same cone before
the register stage. From~the equation tagged eq:witness-linear-shares (in the main paper) we have
$(x^{2})^{(0)} = (x^{(0)})^{2}$, so applying inverse squaring in
$\GF(2^{8})$---i.e., raising to the $128$th power, since
$y^{256}=y$ and $\gcd(2,255)=1$---yields $x^{(0)}$. Similarly,
$(x^{4})^{(1)} = (x^{(1)})^{4}$ and $\gcd(4,255)=1$ imply that
$y\mapsto y^{4}$ is a bijection on $\GF(2^{8})$, so raising
$(x^{4})^{(1)}$ to the $64$th power gives $x^{(1)}$ (since
$4\cdot 64 = 256 \equiv 1 \pmod{255}$). Adding the two recovered
shares yields $x = x^{(0)}\oplus x^{(1)}$ deterministically.
\end{proof}

\subsection{A registered value that is independent of the secret}

The recovery lemma above concerns what a glitch-extended probe sees.
The other half of the gap is a netlist whose \emph{registered} value
carries no information about $x$, so that a verifier inspecting only
register outputs has nothing to flag. We give such a first-order
netlist explicitly; it is the minimal instance of the cone condition
of Theorem~4 (main paper), with $d=1$.

\begin{lemma}
\label{lem:witness-uniform}
Let $x^{(0)},x^{(1)}\in\GF(2^{8})$ be a first-order additive sharing of
$x$ (that is, $x^{(0)}$ uniform and $x^{(1)}=x\oplus x^{(0)}$), and let
$r\in\GF(2^{8})$ be a fresh uniform mask independent of the sharing.
Consider a cone $C$ computing
\begin{equation}
R \;=\; x^{(0)}\oplus x^{(1)}\oplus r ,
\label{eq:witness-R}
\end{equation}
whose output is latched by a register. Then (a) the registered value
$R$ is uniform on $\GF(2^{8})$ and independent of $x$; (b) a single
glitch-extended probe on $C$ observes $x^{(0)}$ and $x^{(1)}$ and hence
recovers $x=x^{(0)}\oplus x^{(1)}$ deterministically.
\end{lemma}

\begin{proof}
(a) Fix $x$ and $x^{(0)}$. Then $x^{(0)}\oplus x^{(1)}=x$ and
$R=x\oplus r$. Since $r$ is uniform and independent of
$(x,x^{(0)})$, $R$ is uniform for every fixed $x$; its distribution
therefore does not depend on $x$. (b) The two shares are inputs of the
same combinational cone feeding the register, so a glitch-extended
probe on $C$ observes both before the register stage; adding them gives
$x$.
\end{proof}

\subsection{The gap}

\begin{theorem}[Concrete gap]
\label{thm:witness-gap}
For the cone $C$ of Lemma~\ref{lem:witness-uniform}:
\begin{enumerate}
\item any value-based verifier that inspects only the distribution of
the registered output $R$ accepts, because $R$ is uniform and
independent of $x$ (Lemma~\ref{lem:witness-uniform}(a));
\item a single glitch-extended probe on $C$ recovers $x$ with
probability one (Lemma~\ref{lem:witness-uniform}(b));
\item the cone-non-completeness lint of Definition~3 (main paper)
rejects, because $C$ contains $d{+}1=2$ distinct shares of $x$.
\end{enumerate}
\end{theorem}

\begin{proof}
Items~1 and~2 are the two parts of Lemma~\ref{lem:witness-uniform}.
Item~3 follows from Definition~3 (main paper): the fan-in of $C$ holds
$x^{(0)}$ and $x^{(1)}$, one more than the $d=1$ shares of $x$ a cone
may contain.
\end{proof}

\begin{remark}[The refresh-free Itoh--Tsujii multiplication]
\label{rem:witness-it}
The cone of $c_{01}$ in the refresh-free design contains
$(x^{2})^{(0)}$ and $(x^{4})^{(1)}$, so by
Lemma~\ref{lem:witness-recover} a glitch-extended probe recovers $x$
and the lint rejects (items~2 and~3 above). We deliberately make no
claim about the marginal of the registered product $c_{01}$ there:
for independent uniform operands the product is not uniform (it equals
$0$ with probability $511/2^{16}$ rather than $2^{-8}$) and its
distribution depends on $x$. The value-blindness part of the gap is
therefore established by the netlist of
Lemma~\ref{lem:witness-uniform}, not by $c_{01}$.
\end{remark}

\section{Pre-barrier PROLEAD history and scope of the second-order instance}
\label{app:mask-reuse}

\textit{History.}  The first-order netlist without the per-cross-term
register barrier (commit \texttt{94623ed}) was evaluated with PROLEAD at
$1{,}536{,}000$ simulations and flagged $160$ of $2{,}816$ probes.  After
the barrier and the operand refresh were inserted, a re-verification at
$512{,}000$ simulations (commits \texttt{d6308ea}, \texttt{6855a66})
reported $0$ of $14{,}400$ flagged probes.  This history is documented in
\texttt{PROLEAD\_\allowbreak{}FINDINGS.md} in the artifact.
\label{app:prolead-findings-historical}

\textit{Independent re-synthesis.}  The gate-level netlist read by the
$512{,}000$-simulation run above was not preserved.  We therefore
re-synthesized the same RTL (same source file) to a different technology
mapping ($46{,}768$ combinational cells and $800$ flip-flops) and re-ran
PROLEAD with the same configuration ($d{=}1$, $19$ cycles, $512{,}000$
simulations, $56$ fresh mask bits).  PROLEAD enumerates $9{,}396$ probing
sets on this netlist, flags $0$ of them, and reports worst
$-\log_{10}(p)=3.61$ against a Bonferroni threshold of $5.27$.  A
gate-level simulation of this netlist reproduces the AES S-box from the XOR
of the two output shares for all $256$ inputs.  The set counts differ from
the earlier run because the mapping differs.

\textit{Scope of the second-order instance.}  The artifact contains a
$d{=}2$ instantiation of the S-box
(\texttt{rtl/masked\_sbox\_second\_order.v}).  We make no security claim
for it.  Its operand refresh XORs one fresh byte into the first two shares
of the $x^4$ operand and passes the third share through unchanged, and an
earlier PROLEAD run on it used a placed probe set whose members each
involve a single share index, so that run could not have exercised a
cross-share probe combination.  Neither the run nor the associated
share-counting argument is part of the evidence of this paper, and both
have been withdrawn from it.

% Bibliography is generated by bibtex from manuscript.bib using the
% IEEEtran BibTeX style (IEEEtran.bst).
% IEEEtran's bibliography is already tightly packed via \IEEEbibitemsep=0pt.


% =====================================================================
% Section added during compression: §Full-Chip Masked AES Round 1
% =====================================================================
\subsection{Full-Chip Masked AES Round~1}
\label{sec:fullchip}

To verify that our masking construction composes correctly at the
chip level, we instantiate 16 parallel first-order masked S-boxes
the first-order masked S-box (\texttt{masked\_sbox\_first\_order}) and wire the
AES-128 round-1 linear layers---ShiftRows, MixColumns, and
AddRoundKey---as share-wise combinational/registered operators
(linear maps are masking-trivial).  The resulting
$\texttt{masked\_\allowbreak{}aes\_\allowbreak{}round1\_\allowbreak{}first\_\allowbreak{}order}$ module accepts
two 128-bit plaintext shares and two 128-bit round-key shares,
plus 128 fresh mask bytes (16 S-boxes $\times$ 8 bytes each =
1024 fresh bits per round); the round key is supplied externally,
i.e., the key schedule is out of scope.

\paragraph{Latency and area.}
The 19-cycle round latency decomposes as 17~cycles for the
S-box pipeline + 1~cycle for the AddRoundKey register +
1~cycle for the output register.  $\texttt{valid\_\allowbreak{}out}$ is a
combinational AND of the 16 S-box valid signals; it pulses for one
cycle when the S-box outputs register (cycle 17), two cycles ahead
of the registered round outputs, so downstream logic must sample
the data path on the registered outputs.  Gate-level
synthesis with Yosys~0.67 ($\texttt{proc} \to \texttt{opt}
\to \texttt{fsm} \to \texttt{memory} \to
\texttt{techmap} \to \texttt{clean}$, no $\texttt{abc}$ in order
to preserve the $\texttt{\$\_\allowbreak{}DFFE\_\allowbreak{}PN0P\_\allowbreak}$/$\texttt{\$\_\allowbreak{}DFF\_\allowbreak{}PN0\_\allowbreak}$
register primitives required for iverilog gate-level
simulation) yields 605{,}698 cells and 13{,}330 D flip-flops
(12{,}544 with enable), with
16 submodule instances of the first-order S-box.

\paragraph{Cell and flip-flop counts.} The $77{,}587$ cells quoted for the
first-order S-box are counted in the PROLEAD cell library, which contains no
multiplexers or other complex gates, so these are decomposed into AND/OR/NOT;
the round-1 total of $605{,}698$ cells is counted in Yosys generic primitives.
The two libraries differ, and $16\times 77{,}587\approx 1.24\times 10^{6}$ is
therefore not a meaningful expectation for the round-1 total; we do not
attribute the difference to individual synthesis effects. The flip-flop count
can be reconciled exactly: the S-box netlist has $784$ flip-flops with enable
and $34$ without ($818$ in total), so $16\times 784=12{,}544$ matches the
number of enable flip-flops of the round, and
$13{,}330 = 16\times 818 + 242$, the remaining $242$ flip-flops belonging to
the round-1 wrapper. Full Yosys statistics logs are released with the
artifact repository for auditability.

\paragraph{Functional verification.}
We generate 100 random $(\text{plaintext}, \text{key})$ pairs
together with 100 sets of 112 random mask bytes and compare the
unmasked output $\texttt{y0\_\allowbreak{}out} \oplus \texttt{y1\_\allowbreak{}out}$ against
a pure-Python reference implementing the same composition
$\mathrm{MC}(\mathrm{SR}(\mathrm{SB}(x))) \oplus \mathit{rk}$
($\texttt{reference.masked\_\allowbreak{}aes}$).  \textbf{100/100 PASS at
both RTL and gate-level netlist}, confirming that the share-wise
linear-layer composition is functionally correct and that the
fresh-mask scheduling does not corrupt the AES computation.

\paragraph{Security.}
Masking correctness at the full-chip level follows from
Theorem~2 (main paper): each of the 16 S-box instances is
cone non-complete and robust $1$-probing secure (the per-S-box
PROLEAD-equivalent statistical analysis in
Section~7 (main paper) returns 0 leaky probes at $d{=}1$), and
ShiftRows, MixColumns, and AddRoundKey are $\mathbb{F}_{2^8}$-linear
maps applied share-wise, so a single probe in the linear layers sees one share index of each
independently shared byte (Theorem~2 (main paper)).  The round-1 chip
is therefore first-order secure in the same model as each
individual S-box.  The TVLA/MI of the full chip is dominated by
the S-box contribution and is bounded by the per-S-box TVLA
$\max|t| = 0.92$ reported in Table~IV (main paper).  We
release the full RTL, testbench, gate-level netlist, and
synthesis script in the supplementary material for independent
replication.


% =====================================================================
% Section added during compression: §10-Round Extension (Skeleton)
% =====================================================================
\subsection*{10-Round Extension (Skeleton)}

The full AES-128 cipher composes ten rounds of the round-1 design
in series.  We provide a skeleton
the full 10-round first-order masked AES skeleton (\texttt{masked\_aes\_10rounds\_first\_order}) that
instantiates ten copies of the existing round-1 module (\texttt{masked\_aes\_round1\_first\_order})
module ten times.  Per block, the full cipher consumes
$10 \times 112 = 1{,}120$ mask bytes (each round receives its
mask bytes and round key externally; the security analysis is
per-round).
The total latency is $10 \times 19 = 190$~cycles per block.

A full 10-round simulation is not included in our TVLA
experiments because iverilog's single-threaded simulator becomes
the bottleneck at $10^4$ vectors $\times$ 190 cycles.  A
Verilator (or commercial VCS / Modelsim) port is the natural
path to scale trace counts; the VCD post-processing
(\texttt{syn/\allowbreak{}vcd\_\allowbreak{}to\_\allowbreak{}hd.py}) is simulator-agnostic and accepts
VCD output from any tool.


% =====================================================================
% Section added during compression: §Discussion paragraphs 2-4 (full)
% =====================================================================
\section{Discussion (full)}
\label{app:discussion-full}

($56$~bits per S-box invocation---six DOM domain-crossing bytes plus
one operand-refresh byte for the first multiplication,
Table~VI (main paper)) is exactly linear in
the number of nonlinear operations, with no hidden refresh network.
This confirms that the random-mask reuse enabled by
Corollary~2 (main paper) is not just a formal but also a
gate-level-visible saving.

Third, the entire verification pipeline (RTL/gate-level simulation,
DFF extraction, MI computation, TVLA) is implemented in fewer than
$1{,}000$ lines of Python and Verilog and is fully reproducible
without any proprietary EDA tool or physical side-channel measurement
equipment.  We release the source alongside the paper.

Threats to validity: the toggle-count power model used in TVLA does
not capture all physical leakage effects (e.g., Hamming-weight bias
of capacitive coupling, power-supply noise); the power testbench
monitors the 34 output-share and valid-path registers
rather than all 818 DFFs of the synthesized netlist, so
switching activity deep inside the S-box pipeline is outside the
model; the PROLEAD-equivalent
verifier is statistical and inherits the finite-sample noise floor of
its MI estimator (we therefore report Miller--Madow-corrected values
and permutation-null percentiles alongside raw plug-in estimates);
first-order HD-mean TVLA has provably low detection power against
transitions averaged over random predecessor state
(Section~7 (main paper)), so the verdicts rest on the conjunction
with the value-based instruments;
the synthesis flow uses Yosys's internal generic
gate library, so absolute cell counts should be interpreted as a
relative ranking rather than as post-place-and-route area numbers.

Empirical scope. Our construction is order-generic in the sense
that it is proven robust $d$-NI with perfect simulation at the gadget level for
arbitrary $d$ (Theorem~1 (main paper)); the
reported PROLEAD, MI, and TVLA numbers instantiate the flow at
$d{=}1$ (post-barrier netlist, $512{,}000$ simulations, $0/14{,}400$
flagged); no empirical claim is made for $d{>}1$.  The cone-non-completeness lint of
Section~5 (main paper), in contrast, is stated and proven at first order: soundness (Theorem~3 (main paper)) and the
witness (Theorem~4 (main paper)); its extension to higher orders is
not claimed here.  A full-DFF power testbench and an
FPGA-based real-silicon pilot are planned as camera-ready extensions and
do not affect the paper's central claims, which are the soundness gap
characterization, the sound-by-construction lint, and the
perfect-simulation proof.

Scope of the lint cross-check. The statistical test used for the lint
table of the main paper is a first-order, fixed-vs-fixed $G$-test on two
secrets with constant inputs and a $32$-bit frontier key; it is evidence for
the lint verdicts, not a proof, and it does not cover cones with more than
$32$ frontier inputs (variant~C). A variant that bypasses the single-domain
pipeline registers is not reported: our re-implementation of that bypass does
not coincide with the one used in earlier drafts, flags $16$ cones, and none of
them leaks in the statistical test.


% =====================================================================
% Moved from main paper (compression): Further automated verifiers
% =====================================================================
\section{Further automated verifiers}
\label{app:further-verifiers}

\paragraph{Further automated verifiers.}
MaskVerif~\cite{maskverif2015} verifies probing-style security of
masked implementations;
IronMask~\cite{ironmask2022} gives algebraic verification of masking
security; and
VERICA~\cite{verica2022} targets combined leakage-and-fault attacks.
These tools are decision procedures (exact or
bounded) over the robust probing model and represent the closest
prior work on automatic verification of masked hardware. Our lint
is complementary: it does \emph{not} decide probing security, but
certifies---by Theorem~3 (main paper)---that a PASS rules
out the structural class of leakage to which the
equipment-free statistical channel is blind, thereby certifying
when the cheaper channel's verdict entails the glitch-extended
one. None of the above tools addresses this cross-channel
soundness question.

Closest to our setting are works that verify concrete masked AES
S-box implementations.  Bloem \emph{et~al.}~\cite{bloem2018formal}
proved the probing security of a full masked first-order AES S-box in
the presence of glitches with exact formal methods, as
SILVER~\cite{knichel2020silver} does;
PROLEAD~\cite{muller2022prolead} trades exactness for
simulation-based statistical coverage of much larger netlists.  On
the implementation side, Cassiers
\emph{et~al.}~\cite{cassiers2021tight} analyzed the probing-security
order of ISW-gadget AES S-boxes, and
Shahmirzadi \emph{et~al.}~\cite{shahmirzadi2021records} set
first-order secure AES performance records verified with SILVER.
These works contribute new gadgets, new implementation records, or
new statistical/model-checking verifiers. Our contribution is
complementary along two axes. First, we introduce a new
\emph{static} verifier---the cone-non-completeness lint of
Corollary~3 (main paper)---that is provably sound with respect to
the robust probing model (Theorem~3 (main paper)) and
provably tight in the sense that any violation admits a concrete
recombination witness (Theorem~4 (main paper),
Section~S2); unlike SILVER's exact but exponential
methods and PROLEAD's simulation-based statistical coverage, the
lint runs in linear time on flattened netlists and returns a
deterministic verdict. Second, we show that an MPC-derived
probing-security proof (Theorem~1 (main paper)), the public PROLEAD toolchain,
and our static lint can mutually cross-check a single fully open
design---including an instructive case in which PROLEAD's
glitch-cone extension detected a deterministic first-order leak
that value-based statistical instruments did not flag
(Section~7 (main paper)), a discrepancy that the new lint
resolves at the netlist level.

The random-probing model (RPM) of~\cite{belaid2020random} bridges
deterministic probing security and noisy-leakage attacks and has spurred
a line of work on tighter composition frameworks: parallel
gadgets~\cite{berti2023parallel}, algebraic RPM bounds
with $O(n^2)$ complexity~\cite{jahandideh2024algebraic}, precomputable
compilers tolerating constant leakage probability~\cite{wang2025precomputation},
and the cardinal-RPC framework
of~\cite{belaid2025cardinal}. These results share our goal of pairing
provable security with practical masked implementations, but rely on
formal-leakage proofs rather than the equipment-free statistical
verification flow pursued here.

Our work differs by coupling an MPC-derived probing-security proof
with a new sound-by-construction static lint and an equipment-free
simulation verification flow, an explicit
randomness--area--latency characterization, and a fully open artifact
in which every reported number is regenerable from the public
repository. The three verification channels---static lint,
robust-probing proof, and PROLEAD/TVLA---agree on the same netlist,
collectively closing the soundness gap identified in
Section~5 (main paper).


% =====================================================================
% Moved from main paper (compression): verification summary
% =====================================================================
\section{Verification metrics (full)}
\label{app:verif-metrics-full}

The PROLEAD-equivalent statistical verifier and the gate-level TVLA on
simulated Hamming-distance power traces return no-detection verdicts
at the designed order.  The full numerical record follows.

statistical verifier: the maximum plug-in MI over all monitored DFF
wires is
$0.022$~bits with a Miller--Madow-corrected value of $\approx 0$
(max 2-wire joint plug-in MI $0.063$~bits excluding the structural
output-share recombination pair, corrected $\approx 0.008$,
both at or below the permutation-null percentile).  Gate-level TVLA on
(corrected) Hamming-distance power traces
found no $|t|{>}4.5$ at any pipeline stage (max $|t| = 0.92$ over
genuine-variance stages; higher-order moments up to $1.28$),
and the bias-corrected mutual information between the simulated power
trace and any of the 8 secret bits is $\approx 0$ bits.  A positive
control on an unmasked S-box confirms detection power (per-wire MI
$1.0$~bit, $43\times$ its null; 2-wire joint MI $2.0$~bits,
$31\times$ its null).

% =====================================================================
% Moved from main paper (compression): §Conclusion 2nd paragraph
% =====================================================================
\section{Full-chip details and future-work directions}
\label{app:conclusion-extra}

We further
instantiated the first-order S-box 16~times to build a full
first-order masked AES-128 round~1
($\texttt{masked\_aes\_round1\_first\_order}$),
consuming $896$ fresh mask bits per round, with a 19-cycle
latency and a 605{,}698-cell / 13{,}330-DFF gate-level
netlist; 100/100 functional vectors pass at both RTL and
post-synthesis.

The formal
full-chip robust $d$-probing security of
Theorem~2 (main paper) follows from cone non-completeness
of the round netlist and the share-wise linearity of ShiftRows, MixColumns,
and AddRoundKey.  By transferring probing-security reasoning from
MPC to the gate level, adding a sound static verifier that closes
the value-based/robust-probing soundness gap, and validating the
result without physical equipment, we narrow the long-standing gap
between claimed and observed side-channel resistance and offer three
mutually consistent verdicts on the same netlist.

Future work extends the
analysis to higher-order instances of the construction, to the
random-probing model of~\cite{belaid2020random},
to on-the-fly mask generation, and to a full 10-round masked
AES-128 accelerator.

\bibliographystyle{IEEEtran}
\bibliography{manuscript}
\end{document}
